\chapter{RQ1 — Global Representation and Multimodal Interaction}
\sloppy

\section{Scope of the evidence}
\label{sec:ch4-scope}
RQ1 asks to what extent lightweight heads over frozen global image and question
embeddings can answer closed-set GQA questions, and how explicit interaction
features, encoder choice and answer-set size affect development accuracy. This
chapter answers that question using development-set evidence only. The main
top-100 view contains 7,714 in-vocabulary questions over 768 represented
images. The top-1000 view contains 9,823 questions over 776 images, while
raw-distribution accuracy uses the shared 10,004-question development
distribution. Results from different denominators are kept separate.

The evidence is ordered by the strength of the comparison. Controlled
low-data contrasts between global heads provide the main answer to RQ1.
Matched-capacity controls then narrow what can be attributed to interaction
features. Visual reliance and training scale describe further properties of
those heads, while the encoder and reasoner comparisons provide bounded
context across less closely matched systems. This order prevents a descriptive
cross-family result from replacing a cleaner controlled comparison.

Unless stated otherwise, values are means over five training seeds. The sample
standard deviation (SD) describes variation between trained seeds, whereas the
image-clustered interval describes evaluation-sampling uncertainty conditional
on those seeds, as defined in Section~3.9. Intervals are nominal and unadjusted
for multiple comparisons. An interval that includes zero does not establish
equivalence or the absence of an effect. No clean-test result exists, so this
chapter makes no confirmatory or held-out performance claim.

\section{Unimodal baselines and multimodal gain}
\label{sec:ch4-baselines}
Table~\ref{tab:ch4-progression} reports the controlled 40k comparison on the
7,714-row top-100 development view. The image-only head reached a mean accuracy
of 0.23438, with seed SD 0.00033 and image-clustered interval [0.22586,
0.24374]. The question-only head was much stronger at 0.45797, with seed SD
0.00290 and interval [0.44717, 0.46917]. The frozen text representation thus
carries substantial task signal, while the global image vector alone is a weak
basis for selecting among the 100 answers.

Combining the two frozen vectors improved the result. Concatenation reached
0.52398, compared with 0.45797 for question only. Their paired difference was
+0.06601 with interval [0.05395, 0.07747], which excludes zero. This comparison
is capacity-confounded because the multimodal head has a different input size,
but it shows that the correct image adds useful information beyond the question
representation. Fusion, which also includes product and absolute-difference
features, reached 0.53840. The capacity controls below determine how much of
that further increase can be attributed to the explicit interactions.

These results separate two questions. The concatenation comparison asks
whether adding the frozen image representation is useful at all. The fusion
comparison asks whether a richer combination of the same global vectors is
associated with a further improvement. Because the head sizes differ, the
second question cannot be answered from this progression alone. The matched
controls in the next section address that limitation directly.

\begin{table}[p]
\centering
\caption[Controlled 40k development accuracies.]{Controlled top-100 development accuracies at \texttt{train\_40k}, using five seeds on 7,714 questions over 768 represented images. Seed SD and the image-clustered interval describe different sources of variation.}
\label{tab:ch4-progression}
\begingroup\scriptsize\setlength{\tabcolsep}{3pt}
\renewenvironment{longtable}[1]{\begin{tabular}{|p{0.28\textwidth}|p{0.16\textwidth}|p{0.14\textwidth}|p{0.25\textwidth}|}}{\end{tabular}}
\let\endhead\relax\renewcommand{\_}{\textunderscore\allowbreak}
% Generated by tools/tab2latex.py; do not edit.
% Presentation metadata: {"header_map":{"95% CI (image-clustered)":"Image-clustered 95% CI","mean accuracy":"Mean accuracy","sd (training seeds, ddof=1)":"Seed SD","system":"System"},"source_headers":["system","mean accuracy","sd (training seeds, ddof=1)","95% CI (image-clustered)"],"thousands":false}
\begin{longtable}{|l|l|l|l|}
\hline
System & Mean accuracy & Seed SD & Image-clustered 95\% CI \\ \hline
\endhead
question only & 0.45797 & 0.00290 & [0.44717, 0.46917] \\ \hline
image only & 0.23438 & 0.00033 & [0.22586, 0.24374] \\ \hline
concat & 0.52398 & 0.00276 & [0.51265, 0.53504] \\ \hline
concat\_wide (matched budget) & 0.53262 & 0.00259 & [0.52154, 0.54394] \\ \hline
fusion\_narrow (matched budget) & 0.53342 & 0.00139 & [0.52130, 0.54546] \\ \hline
fusion & 0.53840 & 0.00218 & [0.52679, 0.55028] \\ \hline
\end{longtable}

\endgroup
\end{table}

\section{Explicit interaction and capacity controls}
\label{sec:ch4-interaction}
The uncontrolled low-data comparison favoured fusion over concatenation by
+0.01442, with seed SD 0.00141 and image-clustered interval [0.00794, 0.02095]
on the same 7,714 development rows. The difference was positive in all five
seeds. This establishes a repeatable difference between the implemented heads,
but it does not isolate the interaction features: fusion has 1,100,388
trainable parameters, compared with 576,100 for concatenation.

The matched controls in Table~\ref{tab:ch4-contrasts} separate feature choice
from capacity. Increasing concatenation to the fusion budget gave +0.00863 over
the original head, with interval [0.00376, 0.01378]. Capacity alone therefore
explains part of the original difference. At this larger budget, fusion
exceeded \texttt{concat\_wide} by +0.00578, but the interval [-0.00116,
0.01226] includes zero, so this contrast is not resolved clearly.

The complementary comparison controls the smaller budget. Reducing fusion to
the concatenation budget still gave +0.00944 over concatenation, with interval
[0.00285, 0.01592]. Full fusion exceeded this narrower version by +0.00498,
with interval [0.00025, 0.00967]. Together, these controls support a bounded
conclusion: capacity explains part, but not all, of the lower-data fusion gain.
They do not show that interaction features are generally preferable to
capacity.

Table~\ref{tab:ch4-contrasts} records both the comparison class and its
uncertainty. This distinction is important because a positive unmatched
contrast and a positive matched contrast support different claims. The former
describes the implemented systems as a whole; the latter controls the
trainable-parameter budget. An interval excluding zero is used as bounded
project evidence, rather than as a general claim about statistical
significance.

Figure~\ref{fig:ch4-ladder} shows the complete low-data ladder. The separate
ablation in VE1-FIG-A1 adds supporting detail. At a matched 576k budget,
product-only and absolute-difference-only heads improved over concatenation by
+0.00957 and +0.00944, while their direct difference was +0.00013 with an
interval crossing zero. Either term recovers much of the matched-budget
improvement, without establishing that the terms implement the same function.

\begin{table}[p]
\centering
\caption[Interaction and capacity contrasts.]{Paired contrasts on the 7,714-row top-100 development view. CAPACITY\_CONFOUNDED comparisons change features and capacity; CLEAN\_PAIRED\_CONTROL comparisons hold the trainable-parameter budget. ``Excludes zero'' is the project evidence criterion, not a significance test.}
\label{tab:ch4-contrasts}
\begingroup\tiny\setlength{\tabcolsep}{1.5pt}
\renewenvironment{longtable}[1]{\begin{tabular}{|p{0.22\textwidth}|p{0.18\textwidth}|p{0.09\textwidth}|p{0.21\textwidth}|p{0.10\textwidth}|p{0.09\textwidth}|}}{\end{tabular}}
\let\endhead\relax\renewcommand{\_}{\textunderscore\allowbreak}
% Generated by tools/tab2latex.py; do not edit.
% Presentation metadata: {"header_map":{"95% CI (image-clustered)":"Image-clustered 95% CI","comparison class":"Comparison class","contrast":"Contrast","effect":"Effect","excludes zero":"Excludes zero","sd (training seeds, ddof=1)":"Seed SD"},"source_headers":["contrast","comparison class","effect","95% CI (image-clustered)","excludes zero","sd (training seeds, ddof=1)"],"thousands":false}
\begin{longtable}{|l|l|l|l|l|l|}
\hline
Contrast & Comparison class & Effect & Image-clustered 95\% CI & Excludes zero & Seed SD \\ \hline
\endhead
concat minus question only & CAPACITY\_CONFOUNDED & 0.06601 & [0.05395, 0.07747] & true & 0.00318 \\ \hline
fusion minus concat & CAPACITY\_CONFOUNDED & 0.01442 & [0.00794, 0.02095] & true & 0.00141 \\ \hline
concat wide minus concat & CLEAN\_PAIRED\_CONTROL & 0.00863 & [0.00376, 0.01378] & true & 0.00302 \\ \hline
fusion minus concat wide & CLEAN\_PAIRED\_CONTROL & 0.00578 & [-0.00116, 0.01226] & false & 0.00298 \\ \hline
fusion minus fusion narrow & CLEAN\_PAIRED\_CONTROL & 0.00498 & [0.00025, 0.00967] & true & 0.00170 \\ \hline
fusion narrow minus concat & CLEAN\_PAIRED\_CONTROL & 0.00944 & [0.00285, 0.01592] & true & 0.00175 \\ \hline
\end{longtable}

\endgroup
\end{table}

\begin{figure}[p]
\centering\includegraphics[width=\textwidth]{ve1/VE1-FIG-02.pdf}
\caption[Low-data ladder with capacity controls.]{Low-data accuracy and paired contrasts on 7,714 development questions over 768 represented images. Intervals are nominal image-clustered evaluation-sampling intervals. The unmatched fusion--concatenation comparison is capacity-confounded; only the matched pairs isolate the feature change, and one matched interval includes zero.}
\label{fig:ch4-ladder}
\end{figure}

\section{Visual reliance}
\label{sec:ch4-reliance}
The visual-reliance experiment changed the image input while leaving each
question unchanged. In v2\_06, image embeddings were either permuted across
development rows or replaced by zeros. The row-level permutation contained 11
self-pairs, or 0.14\% of rows, which slightly dilutes the shuffled
intervention. This is the controlling intervention generation for RQ1.

Figure~\ref{fig:ch4-reliance} shows the drops on the 7,714-row top-100 view.
Fusion dropped by 0.14270 under shuffled images and 0.11880 under zeroed images,
with intervals [0.12769, 0.15759] and [0.10403, 0.13305]. The corresponding
drops for concatenation were 0.11709 [0.10295, 0.13035] and 0.09987 [0.08697,
0.11228]. The matched-budget product head lay between them. The question-only
control had an exact zero drop under both interventions by construction.

The multimodal heads therefore depend on receiving the correct image. This
does not establish visual grounding, compositional reasoning, causal reasoning
or human-like understanding; the intervention measures dependence only.

The two interventions answer related but different questions. Shuffling asks
whether the head needs the image paired with the question rather than another
development image. Zeroing asks whether its prediction changes when the image
representation is removed. Their drop sizes are not treated as a decomposition
of visual information, and neither intervention identifies the reasoning
process used to produce an answer.

\begin{figure}[p]
\centering\includegraphics[width=\textwidth]{ve1/VE1-FIG-05.pdf}
\caption[Visual reliance under wrong and removed images.]{Accuracy drops on 7,714 development questions when the image is replaced or removed. Intervals are nominal image-clustered intervals. Question only is exactly zero by construction and is a control. A drop demonstrates dependence on the correct image, not robust visual reasoning, grounding or compositional understanding.}
\label{fig:ch4-reliance}
\end{figure}

\section{Training scale}
\label{sec:ch4-scaling}
Figure~\ref{fig:ch4-scaling} follows four heads across 40k, 100k and 250k
labelled questions on the 7,714-row top-100 view. Fusion moved from 0.53840 to
0.56308 and 0.58226; concatenation moved from 0.52398 to 0.55084 and 0.57861.
The matched product head reached 0.53355, 0.56365 and 0.58242, while question
only rose from 0.45797 to 0.48294 and 0.49769. Every stored series therefore
increased across the three scales.

The fusion-minus-concatenation difference narrowed from +0.01442 at 40k to a
seed-mean difference of +0.00365 at 250k. Seeds 2 and 42 reversed its direction
at 250k. No image-clustered interval is available for v2\_07, and the figure's
error bars are sample SD across training seeds rather than confidence
intervals. The narrowing is therefore tentative, not statistically
established. Three training sizes also do not support a general scaling law.

The consistent rise in the stored means supports only a result about the three
tested training sizes. It does not show how the heads would behave between or
beyond them. Likewise, the smaller fusion--concatenation gap at the largest
scale is descriptive because its uncertainty is represented only by variation
between training seeds. This limitation is why the scaling claim remains
tentative even though the plotted means are ordered clearly.

\begin{figure}[p]
\centering\includegraphics[width=\textwidth]{ve1/VE1-FIG-03.pdf}
\caption[Training-scale behaviour of the global heads.]{Top-100 development accuracy across 40k, 100k and 250k labelled questions. Error bars are the sample standard deviation across five training seeds (\texttt{ddof=1}); they are not confidence intervals. Under the accepted v2\_07 limitation, no image-clustered interval exists for this family, so the 250k narrowing remains tentative.}
\label{fig:ch4-scaling}
\end{figure}

\section{Representation quality}
\label{sec:ch4-representation}
The primary encoder comparison replaces CLIP with SigLIP while retaining the
40k recipe and the same 7,714 development rows. Figure~\ref{fig:ch4-encoder}
shows SigLIP-minus-CLIP differences of +0.01540 for concatenation, +0.01475 for
fusion and +0.01356 for product. Their intervals are [0.00592, 0.02481],
[0.00601, 0.02401] and [0.00403, 0.02282], all excluding zero. The
question-only difference was +0.00814 with interval [-0.00006, 0.01586], which
includes zero. The stronger multimodal pattern supports representation quality
as an important factor acting mainly through the image path in this experiment.

The comparison is REPRESENTATION\_CONFOUNDED. Encoder identity, embedding
width (512 for CLIP and 768 for SigLIP) and head input width change together.
The development rows match, but training streams are not paired by seed. The
result therefore isolates neither a pure encoder effect nor a single determining constraint.

The cleanest conclusion is therefore limited to the tested encoder systems at
40k. On identical evaluation rows, the multimodal SigLIP systems have positive
contrasts whose intervals exclude zero. The question-only contrast is not
resolved clearly by its interval. This pattern is consistent with an important
image-side representation difference, but the simultaneous width changes stop
it from identifying the encoder alone as the cause.

At 250k on the same 7,714-row view, the stored five-seed SigLIP means were
0.59388 for fusion, 0.59176 for concatenation and 0.50236 for question only,
with seed SDs 0.00317, 0.00261 and 0.00449. No clustered
SigLIP-minus-CLIP interval exists at this scale. As descriptive cross-family
context only, the five-seed SigLIP fusion mean of 0.59388 is comparable to the
three-seed, 21.1M-parameter latent reasoner mean of 0.59580 on that development
view. The capacity-mismatched SigLIP product reference is excluded.

The reasoner comparison is deliberately secondary. It places the stored means
on a common descriptive accuracy scale, but it does not hold architecture,
input granularity, trainable capacity or seed structure fixed. No clustered
contrast connects the two families at this scale. The observation can therefore
show that the values are comparable, but not that one architecture has been
shown to outperform the other.

\begin{figure}[p]
\centering\includegraphics[width=\textwidth]{ve1/VE1-FIG-07.pdf}
\caption[Representation sensitivity at 40k.]{SigLIP-minus-CLIP differences on the same 7,714 development rows. Intervals are nominal image-clustered intervals. Representation is not isolated: encoder identity, embedding width and head input width move together, training streams are not seed-paired, and no clustered contrast exists at 250k.}
\label{fig:ch4-encoder}
\end{figure}

\section{Answer-set size}
\label{sec:ch4-answer-space}
Expanding the vocabulary changes both coverage and the classification problem.
Top-100 coverage is 0.77109 of the shared 10,004-question development
distribution; top-1000 coverage is 0.98191. The top-1000 view therefore has
9,823 questions over 776 images rather than the top-100 view's 7,714 questions
over 768 images.

At 250k, top-1000 in-vocabulary accuracies were 0.40625 for question only,
0.49111 for concatenation, 0.49946 for product and 0.49547 for fusion. Their
intervals in Table~\ref{tab:ch4-answer-space} apply only to the 9,823-row view.
Within that view, fusion minus concatenation was +0.00436 with interval
[-0.00138, 0.00995], so the contrast is not resolved clearly.

The derived raw-distribution column places results on the shared 10,004-row
distribution by multiplying stored in-vocabulary accuracy by coverage. It
gives 0.4490 for top-100 fusion and 0.4462 for top-100 concatenation, compared
with 0.4866 and 0.4822 for their top-1000 counterparts; top-1000 product reaches
0.4905. These are labelled derived quantities, not cross-vocabulary accuracy
contrasts. The in-vocabulary values are never subtracted because answer sets,
training rows and denominators differ. The seed-0 shared-head diagnostic is
omitted because it is not like-for-like with the five-seed top-100 mean.

Coverage and in-vocabulary accuracy answer different questions. Coverage asks
how much of the raw development distribution is represented by an answer
vocabulary. In-vocabulary accuracy asks how often the system is correct after
restricting evaluation to those represented answers. The derived column gives
a row-fair view by placing the stored values back on the shared raw
distribution. It is useful for comparing coverage--accuracy trade-offs, but it
does not turn the two classifier tasks into a paired experiment.

Table~\ref{tab:ch4-answer-space} leaves the question-only derived entry blank
because the controlling report does not state that value. No value is inferred
to fill the presentation. This keeps the table tied to the frozen record and
prevents a convenient arithmetic extension from becoming new evidence.

\begin{table}[p]
\centering
\caption[Answer-space size and development coverage.]{Top-100 and top-1000 results at 250k with distinct in-vocabulary denominators. The raw-distribution column copies the E3 report's derived values on the shared 10,004-row distribution. In-vocabulary accuracies are never differenced across answer spaces.}
\label{tab:ch4-answer-space}
\begingroup\tiny\setlength{\tabcolsep}{1.2pt}
\renewenvironment{longtable}[1]{\begin{tabular}{|p{0.09\textwidth}|p{0.11\textwidth}|p{0.12\textwidth}|p{0.14\textwidth}|p{0.20\textwidth}|p{0.09\textwidth}|p{0.17\textwidth}|}}{\end{tabular}}
\let\endhead\relax\renewcommand{\_}{\textunderscore\allowbreak}
% Generated by tools/tab2latex.py; do not edit.
% Presentation metadata: {"header_map":{"95% CI (image-clustered)":"Image-clustered 95% CI","answer space":"Answer space","coverage":"Coverage","denominator":"Development denominator","derived raw-distribution accuracy":"Derived raw-distribution accuracy","in-vocabulary accuracy":"In-vocabulary accuracy","system":"System"},"source_headers":["answer space","system","denominator","in-vocabulary accuracy","95% CI (image-clustered)","coverage","derived raw-distribution accuracy"],"thousands":false}
\begin{longtable}{|l|l|l|l|l|l|l|}
\hline
Answer space & System & Development denominator & In-vocabulary accuracy & Image-clustered 95\% CI & Coverage & Derived raw-distribution accuracy \\ \hline
\endhead
top-100 & concat & 7,714 & 0.57861 & not available & 0.77109 & 0.4462 \\ \hline
top-100 & fusion & 7,714 & 0.58226 & not available & 0.77109 & 0.4490 \\ \hline
top-1000 & question\_only & 9,823 & 0.40625 & [0.39630, 0.41655] & 0.98191 &  \\ \hline
top-1000 & concat & 9,823 & 0.49111 & [0.47957, 0.50229] & 0.98191 & 0.4822 \\ \hline
top-1000 & product & 9,823 & 0.49946 & [0.48765, 0.51060] & 0.98191 & 0.4905 \\ \hline
top-1000 & fusion & 9,823 & 0.49547 & [0.48321, 0.50775] & 0.98191 & 0.4866 \\ \hline
\end{longtable}

\endgroup
\end{table}

\section{RQ1 summary}
\label{sec:ch4-summary}
On the 7,714-row top-100 development view, lightweight heads over frozen
global representations answer the closed-set task well above either unimodal
reference. At 40k,
explicit interactions improve over concatenation, with an interval excluding
zero and the same direction in every seed. Matched controls narrow the claim:
capacity explains part of the gain, while the lower-budget comparison retains
evidence for the feature change. The multimodal systems also depend on the
correct image.

More labelled questions improve every tested head, although the fusion
advantage narrows at 250k without a clustered interval. The paired 40k encoder
comparison supports representation quality as an important but confounded
factor; at 250k SigLIP fusion is descriptively comparable to the larger
reasoner. The 1,000-answer vocabulary raises coverage and derived
raw-distribution accuracy while defining a different in-vocabulary task.
Chapter~8 considers the broader interpretation and limits of these
development-set findings.

Several stronger conclusions are not established. The evidence does not show
that interaction is generally better than capacity, that representation alone
determines performance, or that the fusion advantage follows a scaling law.
The cross-family context establishes no advantage for SigLIP fusion over the
reasoner, and
the reliance interventions do not establish grounding or compositional
reasoning. Finally, all findings remain development-set results; they provide
no held-out or clean-test conclusion.

\fussy
